\section{Discussion and conclusions}
\label{sec:concl}

\paragraph{What to do.} For a spectral index shared by many pixels with linear amplitudes and known
noise: (i) evaluate the profile likelihood or a proper marginal directly in the acceptance step, never
the $\chi^2$ at freshly drawn amplitudes; (ii) do not use a flat amplitude prior unless the amplitudes
are defined at the pivot frequency of the band set, which depends on the other components and their
priors, and say which reference frequency was used; (iii) quote the profile optimum with the
debiased sandwich error, Eq.~(\ref{eq:debiased}), and check $\Tm/\J_{\rm c}$; (iv) if a prior is put
on a component that does not depend on the index, state the sensitivity of the index to that prior's
variance; (v) model the difference between the beams of the bands, even at $3.7^\circ$ pixels, and
use measured bandpasses for every band including the dust tracer.

\paragraph{What it changes for published numbers.} \pend{Final numbers from Sect.~\ref{sec:data}: the
spread between conventions on the same maps against the spread of published values; the K band
against 30 GHz result.}

\paragraph{What we did not do.} All results are for maps, at low resolution, with one index per
region and Gaussian noise of known covariance. We did not treat time-ordered data, gain or bandpass
sampling, or an index field with inferred smoothness, for which a proper amplitude prior is needed
and has the failure modes of Sect.~\ref{sec:limits}. The real-map test of the tempering law rests on
an acceptance-noise factor measured from the log of the same chain; the forward prediction (widths of
13.1 and 13.6 profile widths with amplitudes stored as maps) has not been run. The LFI noise model is
empirical and isotropic. The misfit of the four-region model on real data is large
(Sect.~\ref{sec:data}), so the quoted statistical errors are lower limits for the regions as defined.

\paragraph{Reproducibility.} Every number in this paper comes from a script in the repository and an
experiment record that states the question, the predictions written before the run, the command and
the outcome, including the predictions that failed. The records are cited inline.
